\chapter{Conclusion}

Infra Monitor combines organization-scoped infrastructure management with
Linux host telemetry, operational event workflows, and machine-learning
features. Its current implementation separates relational operational data
from time-series samples, exposes a Flutter client and Django APIs, and
orchestrates background and model services through the local Compose
configuration (\sourcepath{docker-compose.yml};
\sourcepath{backend/infra_monitor/organization_urls.py}).

The repository contains implementation and test code, but this report does
not assert production deployment, measured availability, classifier accuracy,
or completed user-acceptance testing. These claims require repeatable test
runs and recorded results. The supplied architecture illustration also
omits some current AI paths; this is called out in Chapter 2 rather than
silently treated as complete.

\section*{Future Work}
Priorities for a subsequent iteration include updating the architecture
illustration to show all AI integrations, executing and recording the
documented backend and Flutter checks, adding reproducible performance and
volume tests, and preparing a deployment and monitoring procedure for the
target environment. Semester, submission date, and any project-management
records remain for the team to complete.
